Stein's theorem and its Clifton--Hogarth form are the direct predecessors
\citep{Stein1968,Stein1991,CH1995}. \citet{Rakic1997} takes an event-relative,
transitive ``realized with respect to'' relation, argues that under her
constraints it is not definable from the causal geometry of Minkowski
spacetime, and adds it as a conservative primitive: in our terms, the
added-structure option. Point-based presents \citep{Muller2006}
and the causal-diamond present of Arthur and Savitt
\citep{Arthur2006,Arthur2019,Savitt2009,Savitt2015}, with Dorato's critique
\citep{Dorato2011} and Aames's application to Gödel universes \citep{Aames2022},
are conceptual proposals without a curved-spacetime classification.
\citet{Earman2008} builds worldline becoming in general relativity from causal
pasts and notes that once the symmetries of Minkowski spacetime are lost, the
Clifton--Hogarth definability results no longer apply \citep[p.~151]{Earman2008};
he does not pursue a definability result. In flat
spacetime, definability of sets from relativistic structure is governed by
automorphism groups \citep{MSS2025}.

The category and the naturality schema are those of \citet{BFV2003}; the
locality idea is theirs, and what is offered here is a classification of
relation-valued assignments. The closest structural precedent is the SPASs
property of \citet{FV2012}: theories conforming to a common notion of sameness,
one a subtheory of the other, that are isomorphic in one spacetime should be
isomorphic in all. Theorem~\ref{thm:main} has the same shape, one Minkowski verdict
fixing the alternative everywhere, but its objects are relations between events
and its mechanism is transitivity and (CB). Fewster and Verch show that the full
class of locally covariant theories fails the property, and
Theorem~\ref{thm:weyl} is a counterpart: $\Ichr$ and $\Ichr\cup N$ agree on
$\Mink$ and differ elsewhere. Stein himself posed a cross-model question of a
different kind, whether a cosmological model can be deformed to one with
Gödel-type relations \citep[p.~594]{Stein1970}. Clifton and Hogarth's remark
quoted after Corollary~\ref{cor:worldline} is an observation of difficulty, not a posed
problem. The flat-spacetime counterpart of Corollary~\ref{cor:nopresent}, that
relativity leaves no shared present under its symmetries, is established
\citep{Callender2000,Saunders2002}.

A citation sweep of
\citet{CH1995} and \citet{Stein1991} (Semantic Scholar, 24 October 2026; repeated
with OpenAlex, 25 October 2026), and
reading of the works named here, found no prior extension of
Stein's theorem to curved spacetime. This bounds, but does not certify,
priority.
